% =============================================================================
% 058 / 068
% Status: raw
% =============================================================================
% Notes:
%
% Source question:
% 那 洪武正韻(明朝的官話)又是源自於那裡?
% 有一種謬論是：現在的普通話(北京官話)是受滿州人影響的方言。
% 但其實現代普通話是源自於明朝官話，
% 明朝官話與鮮卑人或吳越人沒有太大的關係？
% =============================================================================

\chapter[那 洪武正韻(明朝的官話)又是源自於那裡? 有一種謬論是：現在的普通話(北京官話)是受滿州人影響的方言。 但其實現代普通話是源自於明朝官話， 明朝官話與鮮卑人或吳越人沒有太大的關係？]{那 洪武正韻(明朝的官話)又是源自於那裡? \\[0.35em] 有一種謬論是：現在的普通話(北京官話)是受滿州人影響的方言。 \\[0.35em] 但其實現代普通話是源自於明朝官話， \\[0.35em] 明朝官話與鮮卑人或吳越人沒有太大的關係？}
\qaanswer{
這得問朱重八為何要將淮西流民區的口語和宋元南京話的音韻體系當作正統了。事實上按『廣韻』、『集韻』的脈絡來看，『洪武正韻』也是不按規矩來的胡鬧。流民軍建立的王朝文化上先天的畸形。明朝官話口語的語法和中古漢語口語出現了斷層，音韻上演化自中古漢語音韻，詞彙還是保留了很多中古漢語詞彙，吸收了各地非官話區和行業行會的詞彙。吳語區因為在南宋城市化達到頂峰，社會多樣性高，市民生活豐富多彩，湧現了大量的戲曲、小說等市民文學作品，創造了大量的詞彙。它們中的很多被明清官話吸收。
}

